\begin{abstract}
Electrochemical impedance spectroscopy (EIS) could enable wider deployment of label-free biosensing, but benchtop instrumentation remains a barrier. Published validations of low-cost alternatives often rely on spectral overlays, error bounds, or significance tests, rather than predefined acceptance criteria. This work presents STEI-ZStat, an AD5941 potentiostat (0.2~Hz--200~kHz) on a 63.1 $\times$ 45.0~mm board, costing US\$25.36 at a build quantity of five. It was validated against a PalmSens4 reference using a Randles dummy cell via a four-stage procedure (reactance-sign screening, resistive plateau checks, CNLS parameter recovery, and equivalence testing). Seven replicates were acquired per instrument. Both reproduced nominal resistance within 0.31\%, with comparable median replicate $|Z|$ standard deviations (0.024\% vs.\ 0.026\%). STEI-ZStat showed lower median magnitude deviation (0.111\% vs.\ 0.236\%). CNLS fitting recovered parameters within $\pm$2\% resistance and $\pm$10\% capacitance limits, with all cross-instrument 95\% confidence intervals within margins. In supplementary measurements, three assembled boards gave inter-board RSDs of 0.027\%, 0.011\%, and 0.148\%. The principal limitation was high-frequency phase error (1.21$^\circ$ at 200~kHz vs.\ 0.33$^\circ$ for PalmSens4). Validation was limited to a passive dummy cell, excluding biological or redox measurements.
\end{abstract}

\begin{IEEEkeywords}
Electrochemical Impedance Spectroscopy (EIS), Analog Front-End (AFE), Potentiostat, Instrument Validation, Biosensor Instrumentation, AD5941
\end{IEEEkeywords}